\newpage
\thispagestyle{empty}

\begin{center}
    \Large\textbf{Abstract}
\end{center}

\vspace{1cm}

Neural networks have revolutionized the field of artificial intelligence, enabling breakthroughs in computer vision, natural language processing, speech recognition, and many other domains. This textbook provides a comprehensive introduction to neural networks, from fundamental mathematical concepts to state-of-the-art architectures and training techniques.

The book is organized into thirteen chapters, covering:

\begin{enumerate}
    \item \textbf{Mathematical Foundations:} Linear algebra (SVD, eigendecomposition, condition numbers), calculus (Taylor expansion, Hessians, convexity), probability theory (MLE, MAP estimation, information theory), and optimization theory.
    
    \item \textbf{Activation Functions:} Comprehensive analysis of sigmoid, tanh, ReLU and its variants (Leaky ReLU, PReLU, ELU), Swish/SiLU, Mish, GELU, and StarReLU, with a proof of the Universal Approximation Theorem and guidance on selection criteria.
    
    \item \textbf{Loss Functions:} Detailed exploration of regression losses (MSE, Huber, MAE), classification losses (cross-entropy, focal loss, label smoothing), and task-specific objectives including contrastive and Dice losses.
    
    \item \textbf{Neural Network Architectures:} Deep Neural Networks (DNN), Convolutional Neural Networks (CNN), and Recurrent Neural Networks (RNN/LSTM/GRU) with mathematical derivations and practical examples.
    
    \item \textbf{Core Operations:} Convolution (standard, dilated, depthwise separable, transposed), pooling, and normalization operations (Batch Norm, Layer Norm, Group Norm, Instance Norm) with mathematical rigor.
    
    \item \textbf{Advanced Mechanisms:} Skip connections (residual, dense), channel attention (SE blocks, CBAM), and spatial attention mechanisms including theoretical analysis of gradient flow.
    
    \item \textbf{Optimization:} Gradient descent variants (SGD, momentum, Nesterov), adaptive optimizers (AdaGrad, RMSProp, Adam, AdamW, AMSGrad), and learning rate schedulers (cosine annealing, cyclical, warm restarts).
    
    \item \textbf{Classic Architectures:} Chronological survey from LeNet-5 (1998) and AlexNet (2012) through VGGNet, GoogLeNet/Inception, ResNet, DenseNet, MobileNet, SENet, EfficientNet, to Vision Transformer (ViT, 2020), U-Net, and ConvNeXt.
    
    \item \textbf{Computational Analysis:} Parameter counting and FLOPs computation for fully connected, convolutional, attention, and normalization layers, with coverage of model compression and scaling laws.
    
    \item \textbf{Training Process and Frameworks:} Practical training procedures, regularization techniques (dropout, weight decay, early stopping), hyperparameter tuning, mixed-precision training, and PyTorch implementation with CUDA concepts.
    
    \item \textbf{Dataset Processing:} Data collection, annotation, augmentation strategies (Mixup, CutMix, AutoAugment, RandAugment), splitting strategies, and data loading pipelines for images, text, and audio.
    
    \item \textbf{Transformer Architecture:} Self-attention mechanism, multi-head attention, positional encodings (sinusoidal, RoPE, ALiBi), encoder-decoder architecture, BERT, GPT, and an overview of modern large language models including LLaMA and scaling laws.
    
    \item \textbf{Modern Training Techniques:} Pre-training and fine-tuning paradigms, parameter-efficient fine-tuning (LoRA, adapters), contrastive learning (SimCLR, InfoNCE), masked autoencoding (MAE), and advanced topics in catastrophic forgetting and multi-task learning.
\end{enumerate}

Each chapter includes detailed mathematical derivations, intuitive explanations, PyTorch code examples, exercises for self-assessment, and references to seminal and recent research papers. Upon completing this textbook, readers will be equipped to design, implement, train, and analyze neural networks for a wide range of real-world applications.

\vspace{2cm}

\textbf{Keywords:} Neural Networks, Deep Learning, Machine Learning, CNN, RNN, LSTM, Transformer, Attention Mechanism, Self-Attention, Optimization, Backpropagation, PyTorch, Computer Vision, Natural Language Processing, Large Language Models, Transfer Learning, Fine-tuning

\newpage

